\chapter{結論と今後の課題}
\label{chap:conclusion}

\section{貢献の概要}

本プロジェクトは、LLMの情動的残差ストリームを解釈可能かつ独立して制御可能な成分へと分解することを目標として設定した。2つの実験を合わせることで、以下の2階層構造が確立される：
\[
\Delta_e = \alpha_R\,\hat{\mathbf{r}}_e + \sum_k \gamma_k\,\mathbf{m}_k,
\]
ここで $\hat{\mathbf{r}}_e$ は感情プロトタイプを選択し、$\mathbf{m}_k$ は感情アイデンティティとは独立して表現的テクスチャを調整する。この分解が本プロジェクトの主要な貢献である。

第1実験では、生のCAAベクトルが共有感情化成分 $\mathbf{g}$ と感情固有のシグナルを混同していることが示され、$\mathbf{g}$ を射影除去することで分類精度とプルチックの円環幾何学の復元性の両方が向上することが明らかになった。第2実験では、感情内残差部分空間が平坦ではないことが示された：プールドPCAは、生成においても因果的に活性であり、かつ感情カテゴリを保持する安定した共有軸を復元する。人間による評価は、知覚的スケールにおいて両知見を裏付けた。

主な貢献は以下の通りである：
\begin{enumerate}
    \item 共有感情化成分を感情固有プロトタイプから分離するCAAベクトルの幾何学的分解。これにより分類精度が向上し、プルチックの円環構造が復元される；
    \item 感情内残差からメタ軸を抽出するプールドPCA手法。$\mathbf{m}_1$ および $\mathbf{m}_2$ は感情カテゴリを変えることなく生成において因果的に活性であり、一方 $\mathbf{m}_3$ は幾何学的に安定しているが、第13層における単層ステアリングでは確実には操作できない；
    \item 分解の両階層が知覚的に識別可能な情動的テクスチャシフトを生成することを確認した人間評価。
\end{enumerate}

\section{結論}

本研究を動機付けた中心的な問いは、LLM出力における情動的テクスチャが感情カテゴリとは独立して分解・制御できるかというものであった。少なくとも本研究で検証した範囲においては、その答えは「はい」である。この分解は幾何学的に原理立てられており、経験的に検証されている。$\hat{\mathbf{r}}_e$ に沿ったステアリングは感情強度をシフトさせ、一方 $\mathbf{m}_k$ に沿ったステアリングは感情ラベルを置換することなく表現的テクスチャをシフトさせる。\par

メタ軸 $\mathbf{m}_k$ は狭い技術的意味において \emph{前言語的} である：すなわち、これらは命名された感情プロトタイプと直交しており、2つの独立した抽出手法から復元可能である。これらが前言語的情動体験の次元に対応するかどうかは、現在の証拠が支持するものの決定的には確立されていない解釈的主張である。モデルには現象的経験はなく、復元されるのは心理学的に根拠のある情動的次元と相関する活性化幾何学である。\par

これらの結果を総合すると、LLMが情動を単に離散的なラベル付き状態の集合としてエンコードしているわけではないことが示唆される。各感情の内部には特定・測定・操作が可能な構造が存在しており、この知見は情動コンピューティング応用とニューラル言語モデルにおける感情的内容のメカニスティックな解釈可能性に対して直接的な含意を持つ。

\section{今後の課題}
\label{sec:future-work}

上述した限界から、いくつかの方向性が直接導かれる。

最も即効性の高い改善策は、LLM-as-judgeスコアを補完するために、メタ軸ステアリングによって生成される情動的テクスチャシフトに対する評価者間の一致度を測定する大規模な人間評価者研究を実施することである。このような研究により、judgeによって評価されたシフトが現在の $N = 18$ パイロットを超えるスケールで知覚的に実在するかどうかが確立される。\par

さらに、介入を複数の層に同時に拡張すること、あるいは抽出されたベクトルを異なるモデルスケールに転移させることにより、分解の頑健性と汎化可能性が検証できる。$\mathbf{m}_3$ が単層ステアリングに応答しないという失敗は、最初のステップとして系統的な多層アブレーションを動機付けている。\par

Jackson ら~\cite{jackson2019emotion} は感情意味論における実質的な異文化間変動を記録している。同一の情動的幾何学が多言語コーパスから復元されるかどうか、またメタ軸ステアリングが言語間で一貫した効果をもたらすかどうかを検証することは、序論~\ref{chap:introduction}で述べた中心的動機の1つに直接対処することになる。\par

プルチックの分類体系を超えて、同一の手法を代替感情モデルに適用することも可能である。異なる感情分類体系（例：RussellのCircumplex VAD次元やデータ駆動型ラベルセット）を用いて構築を繰り返すことで、メタ軸 $\mathbf{m}_k$ がモデル内の安定した情動的構造を反映しているのか、それともプルチック分類体系の人工的産物であるのかを検証できる。\par
